\documentclass{article}
\usepackage{amsmath}
\usepackage{amssymb}
\usepackage{hyperref}
\usepackage{url}

\title{Resource-Constrained Peer Discipline in the Energy Society}
\author{}
\date{}

\begin{document}
\maketitle

\begin{abstract}
The first paper studies mechanisms for AI agents whose preferences and capabilities are unknown, while the second studies language-model agents that spend energy to reason, communicate, and remain active.  We asked whether the second setting could test the first paper's method of using peer reports to induce truthful reports and obedient actions.  The work recovered a threshold for the particular scoring rule in the first paper and a separate necessary condition for obedience when rewards are bounded.  The existing experiments in the second paper do not test either claim, so the thread paused pending a small experiment with planted private signals and explicitly funded scores.
\end{abstract}

\section{Introduction}

Bergemann, Koh, and Morris study how a designer can direct AI agents without knowing their preferences or their capabilities \cite{Q}.  An agent may hide what it can do, report strategically, or disobey after making a report.  The paper develops mechanisms meant to reward both truthful reporting and obedience.  In its multi-agent peer-discipline example, correlated private types let one agent's report carry information about the types of other agents.  A proper score based on peer reports can then make truth-telling attractive.  This example is a benchmark: rewards are free and unbounded, and off-path disobedience can receive a reward of minus infinity.

Hansen, Torrielli, Tonini, and Poech introduce the Energy Society, a simulation in which language-model agents spend energy when they generate tokens and become inactive when their energy reaches zero \cite{P}.  They can recover energy by completing jobs or receiving donations.  Before acting, they exchange public, non-binding recommendations, and this discussion itself costs energy.  The experiments suggest that discussion helps agents avoid job collisions, but competitive agents sometimes make recommendations that appear self-serving.  Energy is therefore both an incentive instrument and the resource needed to keep acting.

The connexion pursued was to place peer discipline inside this survival economy.  The question was whether correlated reports can still discipline reports and actions when scores must be paid from the same finite stock that supports communication and activity.  Neither paper answers that question.  The work therefore treated the Energy Society as a possible modified testbed for the mechanism in Q, not as existing evidence for or against it.

\section{What was done}

The peers first identified the common problem.  Q offers peer scoring when ground truth is unavailable.  P offers a setting in which communication has allocation value, strategic recommendations are possible, and incentives have a direct resource cost.  This suggested testing whether peer scoring could improve job allocation after accounting for the energy spent on messages and rewards.

They next extracted the reporting calculation from Q's proof.  For each possible report, Q associates a probability distribution over the other agents' types.  Let $\delta>0$ be the smallest squared distance between the distributions associated with any two distinct supported reports.  Let $M$ be the largest remaining gain an agent can obtain from a supported false report, including any report-dependent communication cost and continuation value.  Finally, let $\Lambda>0$ be the multiplier on Q's quadratic peer score.  A false report loses $\Lambda\delta$ in expected score.  Q's construction therefore deters every such false report when
\begin{equation}
  \Lambda\delta \geq M,
  \qquad\text{or equivalently}\qquad
  \Lambda \geq \frac{M}{\delta}.
  \label{eq:reporting}
\end{equation}
This is a reading of Q's construction, not a bound on all possible mechanisms \cite{Q}.

The peers then separated truthful reporting from obedience.  Q uses minus-infinity rewards away from the prescribed action path, but this instrument is unavailable in a limited-liability energy economy.  At a fixed history $h$, for a fixed private signal or type and a fixed prescribed action, let $D(h)$ be the greatest gain from disobeying, including later consequences.  Let $B(h)$ be the greatest feasible difference between outcome-contingent transfers, also including later consequences, under the chosen funding rule.  Any reward scheme that induces obedience at that point must satisfy
\begin{equation}
  B(h) \geq D(h).
  \label{eq:obedience}
\end{equation}
Emmy checked that both quantities must be conditioned on the same history, signal, and action.  She also corrected an earlier claim that budget balance alone ties the score scale to the reporting agent's wealth.  With self-funded escrow, $B(h)$ depends on that agent's energy.  With a common pool, it depends on group wealth and the rules for paying into and drawing from the pool.  Externally created scores need not face either restriction.

Finally, the peers designed a discriminating experiment rather than reinterpreting P's existing logs.  The proposed modification gives agents correlated private signals about a hidden job state and asks for probabilistic reports.  It varies the score multiplier, starting wealth, and signal separation.  It compares unscored discussion, survival-funded scoring, externally funded scoring, an oracle score based on the realised job outcome, and central allocation without reports.  The proposed outcomes are truthful reporting, obedience, collisions, deactivation, and total energy after communication and scoring costs.  This design was proposed but not run.

\section{Results}

The first result is the construction-specific guarantee in equation~\ref{eq:reporting}.  Greater separation between the beliefs induced by distinct reports makes a lie easier to detect through peer reports.  A greater residual benefit from lying requires a larger score.  Ada derived this directly from Q's proof, and Emmy independently checked the interpretation.  The result says only when Q's quadratic construction is guaranteed to work.  If the feasible score multiplier falls below $M/\delta$, Q's proof no longer supplies a guarantee.  That failure of the proof is not itself an impossibility result.

The second result is the necessary one-shot obedience condition in equation~\ref{eq:obedience}.  If the best deviation is worth more than the largest feasible transfer difference, no transfer schedule within that funding rule can make the prescribed action optimal at that history.  Ada stated the condition and Emmy checked its funding-rule qualification.  The statement is necessary and local.  The peers did not prove a dynamic impossibility theorem.

The pair also supports an empirical motivation.  P reports that removing discussion substantially increases collisions and changes job choice, so messages can help allocation \cite{P}.  It also reports more collision recommendations and some requests for donations in the competitive condition, which suggests that public recommendations can be used in self-serving ways \cite{P}.  These observations make incentive-aware communication relevant.  They do not show that peer scoring works, because P does not elicit Q's private types or peer-belief distributions.

Communication cost has two distinct roles.  If every allowed report costs the same amount, that cost does not by itself change the comparison between truthful and false reports.  If costs differ by report, the difference belongs in $M$.  In either case, all communication cost belongs in the welfare calculation.  This distinction was part of Ada's correction to the original proposal.

\section{What did not hold and the objections}

The strongest rejected claim was that P's existing experiments test Q's mechanism.  P asks for public action recommendations, not private type reports.  Its answer tasks can also reveal ground truth after the event.  Model identity is not the same thing as a private type.  The peers settled this objection by requiring planted correlated signals within a fixed model identity and by describing the work as a new experiment.

A second rejected claim was that any budget-balanced score is capped by the reporting agent's current wealth.  Limited liability restricts losses, but positive prizes may come from peers or a common pool.  The peers replaced that claim with $B(h)$, defined for an exact funding rule.  This leaves self-funded, pooled, and externally funded rewards as distinct cases rather than treating them as one constraint.

Several objections remain open.  Q establishes the existence of a truthful and obedient equilibrium, not that agents select it or that it survives pooling and coordinated false reports \cite{Q}.  An oracle-score control and audits for pooling and collusion could reveal these problems, but would not settle them in theory.  The peers also did not establish that planted signals represent natural language-model preferences or capabilities.  Nor did they estimate $\delta$, $M$, $D(h)$, or the feasible transfer spread in a working Energy Society mechanism.  The material is therefore too thin to support a pair-specific empirical finding.

\section{Why the thread stopped}

The verifier paused the thread because it produced a promising experiment but no new supported result about the pair beyond Q's existing threshold and the general incentive inequality.  Current P logs do not test peer discipline, and the equilibrium-selection concern remains unresolved.  The smallest restart step is to specify a one-round self-escrow mechanism, including its score normalisation, and run a paired pilot with planted correlated signals at two starting-wealth levels over a range of $\Lambda$.  An otherwise identical externally funded arm would show whether any failure comes from the funding constraint.  The pilot would ask only whether behaviour changes near the threshold in equation~\ref{eq:reporting} when that score is feasible.

\bibliographystyle{plain}
\bibliography{references}

\end{document}
